%!TEX root = TheoretischePhysikC.tex

\section{Vorlesung: Hilbertraum}
Um die Zustände der Quantenmechanik zu beschreiben wird der \emph{Hilbertraum} eine größere Rolle spielen. Um diesen mathematischen Raum zu definieren, benötigen wir erst \emph{komplexe Vektorräume}. Solche Räume werden analog mit dem reellen Fall definiert. 

Warum hier \emph{komplexe} und nicht reelle Vektoren auftreten, wissen wir nicht: die Natur ist nun einmal so gebaut, und eine bessere Alternative wurde bisher nicht gefunden.

Ein \emph{komplexer Vektor} $\mbf{v}\in \mathbb{C}^n$ ist eine Spalte von $n$ komplexen Zahlen
\begin{equation}
	\mbf{v} = 
	\begin{pmatrix}
		v_1 \\ v_2 \\ \vdots \\ v_n
	\end{pmatrix}.
\end{equation}
Es gelten die folgenden Rechenoperationen
\begin{enumerate}
	\item Multiplikation mit Skalar $a\in\mathbb{C}$
	\begin{equation}
		a\mbf{v} = 
	\begin{pmatrix}
		av_1 \\ av_2 \\ \vdots \\ av_n
	\end{pmatrix}.
	\end{equation}
	\item Addition
	\begin{equation}
		\mbf{v}+\mbf{w} = \begin{pmatrix}
		v_1 \\ v_2 \\ \vdots \\ v_n
	\end{pmatrix} + \begin{pmatrix}
		w_1 \\ w_2 \\ \vdots \\ w_n
	\end{pmatrix} =
	\begin{pmatrix}
		v_1+w_1 \\ v_2+w_2 \\ \vdots \\ v_n+w_n
	\end{pmatrix}.
	\end{equation}
	\item Kreuz-Operation
	\begin{equation}
		v^\dag = \begin{pmatrix}
		v_1 \\ v_2 \\ \vdots \\ v_n
	\end{pmatrix}^\dag \equiv \begin{pmatrix}
		v_1^* & v_2^* & \cdots & v_n^*
	\end{pmatrix}
	\end{equation}
	\item Skalarprodukt
	\begin{equation}
		\mbf{v}^\dag\mbf{w} = \begin{pmatrix}
		v_1^* & v_2^* & \cdots & v_n^*
	\end{pmatrix} \begin{pmatrix}
		w_1 \\ w_2 \\ \vdots \\ w_n
	\end{pmatrix} =
	\sum_{j=1}^n v_j^*w_j.
	\end{equation}
\end{enumerate}
Der Raum $V$ aller solcher Vektoren mit festem $n$ wird als \emph{komplexer Vektorraum} bezeichnet. 

Für den Vektorraum selbst genügen die ersten beiden Regeln; Kreuz-Operation und Skalarprodukt sind Zusatzstruktur, die ihn erst zum Hilbertraum macht.

Hierbei wird $n$ als \emph{Dimension} von $V$ definiert. Oft verwendet wird die Notation $\mathbb{C}^n$ für den oben definierten Vektorraum $V$.

\begin{beispiel}
	Der Vektorraum $\mathbb{C}^2$ besteht aus allen zweidimensionalen Vektoren
	\begin{equation}
		\mbf{v} = 
		\begin{pmatrix}
			v_1 \\ v_2
		\end{pmatrix}.
	\end{equation}
	Seien
	\begin{equation}
		\mbf{v} = \begin{pmatrix}
			1+i \\ 1-i
		\end{pmatrix} \quad \text{und} \quad \mbf{w} = \begin{pmatrix}
			1 \\ -1
		\end{pmatrix}
	\end{equation}
	zwei Vektoren. Deren Skalarprodukt lautet
	\begin{equation}
		\mbf{v}^\dag \mbf{w} = \begin{pmatrix}
			1-i & 1+i
		\end{pmatrix} \begin{pmatrix}
			1 \\ -1
		\end{pmatrix} = -2i.
	\end{equation}
\end{beispiel}

Unendlich dimensionale Vektorräume sind ebenfalls möglich und spielen eine sehr wichtige Rolle in der Beschreibung quantenmechanischer Systeme. Der Grund ist physikalisch: unsere besten Modelle sind Felder, und Felder haben unendlich viele Freiheitsgrade. Wie sind solche Vektorräume zu verstehen? Man kann sich unendlich dimensionale Vektoren als Spalten mit unendlich vielen Einträge vorstellen, z.B.:
\begin{equation}
	\mbf{v} = \begin{pmatrix}
		v_1 \\ v_2 \\ \vdots \\
	\end{pmatrix}
\end{equation}
Alle Rechenregeln für Vektorräume gelten auch in den unendlichen dimensionalen Fällen. Das Skalarprodukt $\mbf{v}^\dag \mbf{w}$ wird wie folgt berechnet
\begin{equation}
	\mbf{v}^\dag\mbf{w} =
	\sum_{j=1}^\infty v_j^*w_j.
\end{equation}
Das heißt, in der Regel, dass man in allen Gleichungen $n=\infty$ setzt. Es gibt mehrere subtile Aspekte, die man in Betracht ziehen müsste, um rigorös mit unendlich dimensionalen Vektorräumen zu arbeiten. Allerdings werden wir in dieser Vorlesung diese Aspekte nur nebenbei erwähnen.

Einer davon tritt sofort auf. Der Vektor mit lauter Einsen,
\begin{equation}
	\mbf{v} = \begin{pmatrix}
		1 \\ 1 \\ 1 \\ \vdots
	\end{pmatrix},
\end{equation}
ist eine legale unendliche Liste, aber sein Skalarprodukt mit sich selbst divergiert:
\begin{equation}
	\mbf{v}^\dag\mbf{v} = \sum_{j=1}^\infty 1 = \infty.
\end{equation}
Physikalisch ist die Antwort darauf einfach --- wir lassen solche Vektoren nicht zu. Erlaubt ist ein Vektor genau dann, wenn
\begin{equation}
	\eqbox{\mbf{v}^\dag\mbf{v} = \sum_{j=1}^\infty |v_j|^2 < \infty}
\end{equation}
gilt. Diese Menge ist wieder ein komplexer Vektorraum; man bezeichnet ihn mit $\ell^2$, wobei die Zwei vom Quadrat in der Summe kommt. Damit kann bei der Berechnung von Wahrscheinlichkeiten nichts unendlich groß werden.

Für uns ist das wichtigste Beispiel einen unendlich dimensionalen Vektorraumes der Raum der \emph{quadratintegrierbaren Funktionen} (der als $L^{2}(\mathbb{R})$ bezeichnet wird) welcher $\mathbb{R}$ auf $\mathbb{C}$ abbildet. Hierbei bedeutet quadratintegrierbar, dass
\begin{equation}
	\int_{-\infty}^\infty |f(x)|^2 \, dx < \infty.
\end{equation}
Quadratintegrierbare Funktionen können addiert, subtrahiert, und mit Skalargrößen multipliziert werden. Bezüglich dieser Operationen bleiben die Ergebnisse immer quadratintegrierbar. Daher bildet die Menge aller quadratintegrierbaren Funktionen einen abstrakten Vektorraum. Man betrachtet eine quadratintegrierbare Funktion $f(x)$ als einen Vektor. Auch hier gibt es ein Skalarprodukt; an die Stelle der Summe tritt ein Integral:
\begin{equation}
	\langle f, g\rangle \equiv \int_{-\infty}^\infty f^*(x)\, g(x)\, dx.
\end{equation}

\begin{beispiel}
	Ein \emph{Nicht}-Beispiel von einer quadratintegrierbaren Funktion wird von der konstanten Funktion $f(x) = 1$ geliefert:
	\begin{equation}
		\int_{-\infty}^\infty |f(x)|^2\, dx = \int_{-\infty}^\infty 1\, dx = \infty.
	\end{equation}
	Das ist das Gegenstück zum Vektor mit lauter Einsen von vorhin.
\end{beispiel}

% fig_vl02_raeume.tex — Der Zulassungstest für Hilbertraum-Vektoren:
% endliche Norm.  Links die Summe (ℓ²), rechts das Integral (L²(R)) —
% dieselbe Frage, einmal über einen diskreten und einmal über einen
% kontinuierlichen Index.
%
% Was die Abbildung zeigt, das der Text nicht sagt: es genügt nicht, dass
% die Einträge (bzw. die Funktionswerte) verschwinden — es kommt auf die
% Rate an.  1/j und 1/sqrt(j) gehen beide gegen null; nur die erste Folge
% ist in ℓ².  Das ist ein Skalierungsverhalten, und ein Skalierungs-
% verhalten sieht man, statt es zu erklären.
%
% Geometrie — jede Kurve ist die geschlossene Form der jeweiligen
% Partialsumme bzw. des Partialintegrals, nicht freihand:
%
%   (a)  S_N = sum_{j<=N} |v_j|^2
%          v_j = 1/j      : S_N -> pi^2/6 = 1.6449...   (Basler Reihe)
%          v_j = 1/sqrt(j): S_N = H_N ~ ln N + gamma    (log-divergent)
%          v_j = 1        : S_N = N                     (linear divergent)
%        Die Punkte sind die exakt aufsummierten Werte (N = 1..25); die
%        Asymptote liegt bei pi^2/6, nicht bei einem gerundeten Strich.
%
%   (b)  F(X) = int_{-X}^{X} |f|^2 dx
%          f = e^{-x^2/2}    : F = sqrt(pi) erf(X) -> sqrt(pi) = 1.7725...
%          f = 1/sqrt(1+|x|) : F = 2 ln(1+X)             (log-divergent)
%          f = 1             : F = 2X                    (linear divergent)
%
% Beide Panels teilen die y-Skala (0 bis 4.6), damit die beiden endlichen
% Grenzwerte pi^2/6 und sqrt(pi) direkt vergleichbar nebeneinander liegen
% und die beiden log-Divergenzen erkennbar dieselbe Form haben.
%
% Entwurfsentscheidungen:
%   * Farbe trägt genau eine Bedeutung, nämlich das Ergebnis des Tests:
%     kalt (cgblue) = zugelassen, warm = durchgefallen.  Innerhalb von
%     „warm“ unterscheidet der Ton den subtilen vom groben Fall
%     (birrenorange: Einträge gehen gegen null und trotzdem divergent;
%     birrenred: das Beispiel der Vorlesung, v_j = 1 bzw. f = 1).
%     Zusätzlich unterscheiden Strichmuster die drei Kurven, damit die
%     Aussage den Graustufendruck überlebt.
%   * Die Grenzwerte pi^2/6 und sqrt(pi) sind neutral graue, gestrichelte
%     Referenzlinien (levellinegrey) — sie sind keine Daten.
%   * Panel (a) trägt Marker auf jedem ganzzahligen N, Panel (b) nicht:
%     die Textur der Kurve zeigt selbst, dass links über einen diskreten
%     und rechts über einen kontinuierlichen Index summiert wird.
%   * Die divergenten Kurven laufen absichtlich oben aus dem Rahmen; sie
%     sind dort beschriftet, wo sie ihn verlassen.  Die log-divergente
%     Kurve tut das nicht — sie biegt nur ab und koennte als saettigend
%     missverstanden werden.  Deshalb steht ihr Wachstumsgesetz
%     (~ ln N bzw. ~ 2 ln X) als dritte Zeile an ihrer Beschriftung.
%   * ymax = 5 statt 4.6: die Vorlesung setzt 12pt, dort werden
%     \footnotesize-Knoten spuerbar groesser als in der 10pt-Vorschau.
%     Die Beschriftungen wurden gegen eine 12pt-Vorschau gesetzt, nicht
%     gegen scripts/tikz-preview.sh allein — bei 10pt sah eine Kollision
%     zwischen der blauen Kurvenbeschriftung und der C^n-Randnotiz noch
%     harmlos aus, im Satz ueberlappten die beiden.
%   * C^n erscheint als graue Randnotiz statt als vierte Kurve: dort hat
%     die Summe n Terme und ist immer endlich — es gibt nichts zu prüfen.
%     Das ist der ganze Unterschied zwischen C^2 und den beiden anderen
%     Räumen der Vorlesung.

\begin{center}
\begin{tikzpicture}

% ─────────────────────────────────────────────────────────────
%  (a)  Partialsummen: der Test in ℓ²
% ─────────────────────────────────────────────────────────────
\begin{axis}[
    name=seq,
    tpcplot,
    width=6.6cm, height=5.9cm,
    xlabel={$N$},
    ylabel={$\sum_{j=1}^{N}|v_j|^{2}$},
    every axis y label/.style={at={(axis description cs:-0.14,0.5)},
                               rotate=90, anchor=south, font=\small},
    xmin=0, xmax=26, ymin=0, ymax=5.0,
    xtick={0,5,10,15,20,25}, ytick={0,1,2,3,4},
    grid=major,
    title={(a) \tpcifenglish{Partial sums in $\ell^2$}%
                            {Partialsummen in $\ell^2$}},
]

  % Grenzwert der Basler Reihe: neutrale Referenzlinie, keine Daten
  \addplot[levellinegrey, domain=0:26, samples=2] {1.6449341};

  % v_j = 1/j : S_N = sum 1/j^2 -> pi^2/6
  \addplot[cgblue, line width=1.0pt, mark=*, mark size=1.0pt,
           mark options={draw=cgblue, fill=cgblue}] coordinates {
    (1,1.0000)(2,1.2500)(3,1.3611)(4,1.4236)(5,1.4636)(6,1.4914)
    (7,1.5118)(8,1.5274)(9,1.5398)(10,1.5498)(11,1.5580)(12,1.5650)
    (13,1.5709)(14,1.5760)(15,1.5804)(16,1.5843)(17,1.5878)(18,1.5909)
    (19,1.5937)(20,1.5962)(21,1.5984)(22,1.6005)(23,1.6024)(24,1.6041)
    (25,1.6057)};

  % v_j = 1/sqrt(j) : S_N = H_N, logarithmisch divergent
  \addplot[birrenorange, line width=1.0pt, densely dashed,
           mark=*, mark size=1.0pt,
           mark options={solid, draw=birrenorange, fill=birrenorange}]
  coordinates {
    (1,1.0000)(2,1.5000)(3,1.8333)(4,2.0833)(5,2.2833)(6,2.4500)
    (7,2.5929)(8,2.7179)(9,2.8290)(10,2.9290)(11,3.0199)(12,3.1032)
    (13,3.1801)(14,3.2516)(15,3.3182)(16,3.3807)(17,3.4396)(18,3.4951)
    (19,3.5477)(20,3.5977)(21,3.6454)(22,3.6908)(23,3.7343)(24,3.7760)
    (25,3.8160)};

  % v_j = 1 : S_N = N, das Nicht-Beispiel der Vorlesung
  \addplot[birrenred, line width=1.0pt, dash pattern=on 4pt off 2pt on 1pt off 2pt,
           mark=*, mark size=1.0pt,
           mark options={solid, draw=birrenred, fill=birrenred}]
  coordinates {(1,1)(2,2)(3,3)(4,4)(5.3,5.3)};

  % Beschriftung direkt an den Kurven
  \node[lbl, text=cgblue, align=left, anchor=west]
    at (axis cs:5.4,0.80) {$v_j=\frac{1}{j}$\\[1pt]$\in\ell^2$};
  \node[lbl, text=birrenorange, align=left, anchor=west]
    at (axis cs:14.6,2.45) {$v_j=\frac{1}{\sqrt{j}}$\\[1pt]$\notin\ell^2$\\[1pt]
                          {\scriptsize$\sim\ln N$}};
  \node[lbl, text=birrenred, align=left, anchor=west]
    at (axis cs:5.4,4.35) {$v_j=1$\\[1pt]$\notin\ell^2$};

  % Grenzwert der Basler Reihe, direkt über seiner Referenzlinie
  \node[lbl, text=birrengray, font=\scriptsize, anchor=west]
    at (axis cs:6.0,1.80) {$\pi^{2}/6$};

  % C^n: der endlichdimensionale Fall, in dem nichts zu prüfen ist
  \node[lbl, text=birrengray, font=\scriptsize, align=right,
        anchor=south east]
    at (axis cs:25.6,0.10)
    {$\C^{n}$: $n$ \tpcifenglish{terms}{Terme}\\[1pt]
     \tpcifenglish{always finite}{stets endlich}};
\end{axis}

% ─────────────────────────────────────────────────────────────
%  (b)  Partialintegrale: derselbe Test in L^2(R)
% ─────────────────────────────────────────────────────────────
\begin{axis}[
    tpcplot,
    at={(seq.south east)}, anchor=south west, xshift=11mm,
    width=6.6cm, height=5.9cm,
    xlabel={$X$},
    ylabel={$\int_{-X}^{X}|f|^{2}\,dx$},
    every axis y label/.style={at={(axis description cs:-0.14,0.5)},
                               rotate=90, anchor=south, font=\small},
    xmin=0, xmax=4.3, ymin=0, ymax=5.0,
    xtick={0,1,2,3,4}, ytick={0,1,2,3,4},
    grid=major,
    title={(b) \tpcifenglish{Partial integrals in $L^2(\R)$}%
                            {Partialintegrale in $L^2(\R)$}},
]

  % Grenzwert des Gauss-Integrals: neutrale Referenzlinie
  \addplot[levellinegrey, domain=0:4.3, samples=2] {1.7724539};

  % f = exp(-x^2/2) : F(X) = sqrt(pi) erf(X) -> sqrt(pi)
  \addplot[cgblue, line width=1.0pt] coordinates {
    (0.00,0.0000)(0.10,0.1993)(0.20,0.3947)(0.30,0.5825)(0.40,0.7593)
    (0.50,0.9226)(0.60,1.0703)(0.70,1.2014)(0.80,1.3153)(0.90,1.4125)
    (1.00,1.4936)(1.10,1.5601)(1.20,1.6135)(1.30,1.6555)(1.40,1.6879)
    (1.50,1.7124)(1.60,1.7305)(1.70,1.7437)(1.80,1.7531)(1.90,1.7597)
    (2.00,1.7642)(2.10,1.7672)(2.20,1.7692)(2.30,1.7704)(2.40,1.7712)
    (2.50,1.7717)(2.60,1.7720)(2.70,1.7722)(2.80,1.7723)(2.90,1.7724)
    (3.00,1.7724)(3.20,1.7724)(3.40,1.7725)(3.60,1.7725)(3.80,1.7725)
    (4.00,1.7725)(4.30,1.7725)};

  % f = 1/sqrt(1+|x|) : F(X) = 2 ln(1+X), logarithmisch divergent
  \addplot[birrenorange, line width=1.0pt, densely dashed,
           domain=0:4.3, samples=120] {2*ln(1+x)};

  % f = 1 : F(X) = 2X, das Nicht-Beispiel der Vorlesung
  \addplot[birrenred, line width=1.0pt,
           dash pattern=on 4pt off 2pt on 1pt off 2pt,
           domain=0:2.7, samples=2] {2*x};

  % Beschriftung direkt an den Kurven
  \node[lbl, text=cgblue, align=left, anchor=west]
    at (axis cs:1.60,0.92) {$f=e^{-x^{2}/2}$\\[1pt]$\in L^2(\R)$};
  \node[lbl, text=birrenorange, align=left, anchor=west]
    at (axis cs:2.42,4.05) {$f=\frac{1}{\sqrt{1+|x|}}$\\[1pt]$\notin L^2(\R)$\\[1pt]
                          {\scriptsize$\sim 2\ln X$}};
  \node[lbl, text=birrenred, align=left, anchor=west]
    at (axis cs:0.10,4.35) {$f=1$\\[1pt]$\notin L^2(\R)$};

  % Grenzwert des Gauss-Integrals, direkt über seiner Referenzlinie
  \node[lbl, text=birrengray, font=\scriptsize, anchor=south]
    at (axis cs:3.45,1.82) {$\sqrt{\pi}$};

\end{axis}

\end{tikzpicture}
\end{center}


Achten Sie auf die mittlere Kurve in beiden Bildern: bei $v_j = 1/\sqrt{j}$ und
$f(x) = 1/\sqrt{1+|x|}$ gehen die Einträge sehr wohl gegen null, die Norm aber
logarithmisch gegen unendlich. Nicht das Verschwinden entscheidet, sondern die
Rate.

\begin{greybox}[Drei Räume, drei Systeme]
	Damit haben wir praktisch alle Hilberträume beisammen, die wir brauchen werden: $\mathbb{C}^2$ gehört zum \emph{Qubit}, $\ell^2$ zum \emph{harmonischen Oszillator} (etwa einem Photon in einer Glasfaser) und $L^2(\mathbb{R})$ zu einem \emph{Teilchen in einer Dimension}.
\end{greybox}


